\documentclass[tikz,border=8pt]{standalone}
\usepackage{ctex,amsmath,amssymb}
\usetikzlibrary{arrows.meta,positioning,calc,shapes.geometric}
% Shared visual language for LFS architecture diagrams.
\RequirePackage{../../mymath}

\definecolor{LFSBlue}{HTML}{2563EB}
\definecolor{LFSBlueSoft}{HTML}{DBEAFE}
\definecolor{LFSGreen}{HTML}{15803D}
\definecolor{LFSGreenSoft}{HTML}{DCFCE7}
\definecolor{LFSOrange}{HTML}{C2410C}
\definecolor{LFSOrangeSoft}{HTML}{FFEDD5}
\definecolor{LFSRed}{HTML}{B91C1C}
\definecolor{LFSRedSoft}{HTML}{FEE2E2}
\definecolor{LFSPurple}{HTML}{7E22CE}
\definecolor{LFSPurpleSoft}{HTML}{F3E8FF}
\definecolor{LFSGray}{HTML}{475569}
\definecolor{LFSGraySoft}{HTML}{F1F5F9}

\tikzset{
  lfs picture/.style={
    font=\small,
    node distance=8mm and 12mm,
    >=Latex,
    every path/.style={line cap=round, line join=round}
  },
  block/.style={
    draw=LFSBlue,
    fill=LFSBlueSoft,
    rounded corners=2pt,
    line width=0.8pt,
    align=center,
    inner xsep=7pt,
    inner ysep=5pt,
    minimum height=10mm
  },
  compute/.style={
    block,
    draw=LFSPurple,
    fill=LFSPurpleSoft
  },
  storage/.style={
    block,
    draw=LFSGreen,
    fill=LFSGreenSoft
  },
  interface/.style={
    block,
    draw=LFSOrange,
    fill=LFSOrangeSoft
  },
  risk/.style={
    block,
    draw=LFSRed,
    fill=LFSRedSoft
  },
  neutral/.style={
    block,
    draw=LFSGray,
    fill=LFSGraySoft
  },
  decision/.style={
    diamond,
    draw=LFSOrange,
    fill=LFSOrangeSoft,
    line width=0.8pt,
    align=center,
    inner sep=2pt,
    aspect=2
  },
  group/.style={
    draw=LFSGray,
    rounded corners=3pt,
    dashed,
    line width=0.7pt,
    inner sep=6pt
  },
  arrow/.style={
    -{Latex[length=2.4mm,width=1.6mm]},
    draw=LFSGray,
    line width=0.9pt
  },
  data arrow/.style={
    arrow,
    draw=LFSBlue
  },
  control arrow/.style={
    arrow,
    draw=LFSOrange,
    dashed
  },
  residual/.style={
    arrow,
    draw=LFSGreen,
    rounded corners=4pt
  },
  label text/.style={
    font=\scriptsize,
    text=LFSGray,
    fill=white,
    inner sep=1.5pt
  },
  title/.style={
    font=\bfseries,
    text=LFSGray,
    align=center
  }
}


\begin{document}
\begin{tikzpicture}[lfs picture]

% 左侧：经典信息熵维恩图
\begin{scope}[shift={(0,0)}]
  % 填充左右两圆
  \fill[LFSBlueSoft, opacity=0.85] (0,0) circle (1.8cm);
  \fill[LFSGreenSoft, opacity=0.85] (2.2,0) circle (1.8cm);

  % 描边
  \draw[LFSBlue, line width=1.1pt] (0,0) circle (1.8cm);
  \draw[LFSGreen, line width=1.1pt] (2.2,0) circle (1.8cm);

  % 文字标签
  \node at (-0.8,0) [align=center] {\bfseries $H(X\mid Y)$\\{\scriptsize 条件熵}};
  \node at (1.1,0) [align=center] {\bfseries $I(X;Y)$\\{\scriptsize 互信息}};
  \node at (3.0,0) [align=center] {\bfseries $H(Y\mid X)$\\{\scriptsize 条件熵}};

  % 上方与下方标注
  \node at (0, 2.2) [text=LFSBlue, font=\small\bfseries] {$H(X)$ 边缘熵};
  \node at (2.2, 2.2) [text=LFSGreen, font=\small\bfseries] {$H(Y)$ 边缘熵};
  \node at (1.1,-2.3) [text=LFSGray, font=\small\bfseries] {$H(X,Y)$ 联合熵（整体区域）};
\end{scope}

% 右侧：交叉熵与相对熵（KL散度）分解条形图
\begin{scope}[shift={(6.2,-1.5)}]
  \node[title] at (2.4, 3.7) {交叉熵与相对熵分解关系};

  % 完整交叉熵条形外框
  \draw[draw=LFSOrange, fill=LFSOrangeSoft, line width=0.9pt, rounded corners=2pt] (0,1.8) rectangle (4.8, 2.8);
  \node at (2.4, 2.3) {\bfseries $H(p, q) = -\sum_x p(x) \log q(x)$ \quad（总交叉熵损失）};

  % 分解的两段
  \draw[draw=LFSBlue, fill=LFSBlueSoft, line width=0.9pt, rounded corners=2pt] (0,0.2) rectangle (2.1, 1.2);
  \node at (1.05, 0.7) [align=center] {\bfseries $H(p)$\\{\scriptsize 真实分布固有熵}};

  \draw[draw=LFSRed, fill=LFSRedSoft, line width=0.9pt, rounded corners=2pt] (2.3,0.2) rectangle (4.8, 1.2);
  \node at (3.55, 0.7) [align=center] {\bfseries $D_{\mathrm{KL}}(p\|q)$\\{\scriptsize 预测偏差相对熵}};

  \node at (2.2, 0.7) [font=\large\bfseries] {$+$};

  % 关系指示虚线
  \draw[control arrow] (1.05, 1.2) -- (1.05, 1.8);
  \draw[control arrow] (3.55, 1.2) -- (3.55, 1.8);

  \node[label text] at (2.4, -0.3) {优化目标：$\min_q H(p,q) \iff \min_q D_{\mathrm{KL}}(p\|q)$};
\end{scope}

\end{tikzpicture}
\end{document}
